\subsection{Function Requirements}

Functions shall:

\begin{itemize}
	\item have one primary responsibility;
	\item have a clearly defined contract;
	\item validate input parameters;
	\item return an explicit status for recoverable failure;
	\item avoid hidden side effects;
	\item avoid modifying caller data unless documented;
	\item remain short enough for complete review.
	\item Limit nesting depth to four levels.
\end{itemize}

Functions should normally have no more than 10 local variables. If more are required, the function should be refactored into smaller functions.

Every function that may be accessed from outside its module MUST include a function prototype (declaration).

Function declarations and definitions that accept no parameters MUST explicitly use (void), never empty parentheses ().

Function implementations MUST place the return type and any storage-class or qualifier keywords on a separate line from the function name.

A function SHOULD have at most 2 return points: one early return at the top of the function for parameter/argument validation, and one return at the end of the function. Avoid return statements in the middle of the function body; use a return-value variable to accumulate the result instead.



\subsection{Parameter Validation}

Public functions shall validate all pointer parameters, lengths, ranges, and state assumptions that are not guaranteed by the type system.

\subsection{Return Values}

The caller shall check every non-void return value.

If a return value is intentionally ignored, the reason shall be explicit.

\begin{lstlisting}[language=C]
(void)log_message("startup complete");
\end{lstlisting}
